% Copyright (c) 2026 Eugene Joseph M. Ragasa.
% SPDX-License-Identifier: Apache-2.0
% Author: Eugene Joseph M. Ragasa
% Formal proof completed: 2026-10-20
\documentclass[11pt]{article}
\usepackage[T1]{fontenc}
\usepackage{amsmath,amssymb,booktabs,tabularx}
\usepackage[margin=1in]{geometry}
\usepackage{xurl}
\usepackage[hidelinks]{hyperref}
\setlength{\emergencystretch}{3em}
\usepackage{listings}
\lstset{basicstyle=\ttfamily\small,breaklines=true,columns=fullflexible,
  literate={ℝ}{{$\mathbb R$}}1 {ℂ}{{$\mathbb C$}}1 {•}{{$\mathbin\bullet$}}1}
\newcommand{\StatementID}{BG-PROJECTION-THM-001}
\newcommand{\LeanDeclaration}{ProjectKoios.BandGauge.Projection.twoBandParentLaurentMatrix\_sq}
\title{Project Koios Formal Proof Dossier\\\large The Rank-Two Laurent Parent Squares to Its Scalar Radicand}
\author{Eugene Joseph M. Ragasa}
\date{Formal proof completed: 2026-10-20}
\begin{document}
\sloppy
\maketitle

\section{Identity and status}
\begin{tabularx}{\textwidth}{@{}lX@{}}\toprule
Statement ID & \path{BG-PROJECTION-THM-001} \\
Status & proved; no admissions \\
Lean declaration & \path{ProjectKoios.BandGauge.Projection.twoBandParentLaurentMatrix_sq} \\
Owning module & \path{ProjectKoios/BandGauge/Projection/LaurentObstruction.lean} \\
Author and license & Eugene Joseph M. Ragasa; \texttt{Apache-2.0} \\
Formal proof completed & 2026-10-20 \\
\bottomrule\end{tabularx}

This status describes one algebraic Lean theorem. It is not a claim that the full
spectral-projection proposition, any numerical result, or the publication is formally
complete.

\section{Exact mathematical statement}
Let $\Delta,t\in\mathbb R$ and work in the complex Laurent-polynomial ring
$\mathbb C[z,z^{-1}]$. Define
\[
 H(z)=\begin{pmatrix}
 \Delta&t(1+z^{-1})\\
 t(1+z)&-\Delta
 \end{pmatrix},\qquad
 r(z)=\Delta^2+2t^2+t^2z+t^2z^{-1}.
\]
Then
\[
 H(z)^2=r(z)I_2.
\]
There is no nondegeneracy premise on this supporting identity.

\section{Source provenance}
\begin{tabularx}{\textwidth}{@{}lX@{}}\toprule
Owner repository & \texttt{ksdft2effmass} \\
Navigation branch & \texttt{work/paper-3-higher-rank-extension} \\
Immutable commit & \path{99a46cb59d0c6f27ada7d709f194996ef9619d9f} \\
Source path & \path{docs/publications/papers/Ragasa2026b/v2/manuscript.tex} \\
Source locator & lines 316--330, especially the matrix and radicand defining the eigenvalues \\
Source SHA-256 & \path{d8d5951d8d7ed75355b5d94f37fcbddc70c1466897a5770dbe9f877e60016daf} \\
Dirty-state limitation & cited bytes are the clean immutable commit recorded by the owner-source audit \\
\bottomrule\end{tabularx}

The Lean theorem \emph{adapts} the source matrix into a Laurent-polynomial identity. It
is supporting algebra for, but is not identical to, the source proposition.

\section{Assumptions, conventions, and statement delta}
\begin{tabularx}{\textwidth}{@{}XXX@{}}\toprule
Source convention & Lean encoding & Delta or limitation \\
\midrule
One momentum variable $z=e^{ik}$ & integer Laurent exponents in
\texttt{LaurentPolynomial $\mathbb C$} & unit-circle evaluation is not used by this algebraic declaration \\
Rank two in the displayed basis & \texttt{Matrix (Fin 2) (Fin 2)} & basis is fixed exactly \\
Real $\Delta,t$ & exact coercion into $\mathbb C$ & no floating-point data \\
Source proposition assumes $\Delta t\ne0$ & this algebraic identity is proved for all
$\Delta,t$ & broader validity is used only as standard supporting algebra, not as a
broader projection claim \\
Nearest-neighbor support & only exponents $-1,0,1$ occur in the definition & no physical
metric or hopping-distance theorem \\
\bottomrule\end{tabularx}

\section{Statement-specific prior art}
\begin{tabularx}{\textwidth}{@{}lXXX@{}}\toprule
Relation & Canonical identity & Locator & Role and caveat \\
\midrule
adapts & owner source above; canonical literature identity unresolved &
\texttt{manuscript.tex:316--330} & supplies the exact parent and radicand; no priority
claim is made for matrix multiplication \\
\bottomrule\end{tabularx}

No accepted external reference is asserted for this elementary identity. A missing
reference record would not make the algebra novel.

\section{Motivation and contribution boundary}
The identity exposes why the scalar square root in the two-band eigenvalues is precisely
the Laurent radicand used by the obstruction theorem. Its value is traceable exact
bookkeeping between the parent and the algebraic obstruction. Matrix multiplication and
Laurent monomial arithmetic are standard; this dossier makes no novelty claim for them.

\section{Proof schematic}
\[
\boxed{H(z),\ r(z)}
\Longrightarrow
\boxed{z^{-1}z=zz^{-1}=1}
\Longrightarrow
\begin{cases}
(H^2)_{12}=(H^2)_{21}=0,\\
(H^2)_{11}=(H^2)_{22}=r(z)
\end{cases}
\Longrightarrow
\boxed{H(z)^2=r(z)I_2}.
\]

\section{Human-readable mathematical argument}
Multiply the two matrices. The upper-right entry is
$\Delta t(1+z^{-1})-\Delta t(1+z^{-1})=0$, and the lower-left entry cancels in the same
way. For the first diagonal entry,
\[
 \Delta^2+t^2(1+z^{-1})(1+z)
 =\Delta^2+t^2(2+z+z^{-1})=r(z),
\]
because $z^{-1}z=1$ in the Laurent ring. The other diagonal entry has the factors in
the opposite order and gives the same value. Thus the square is $r(z)I_2$.

\section{Lean realization}
\begin{lstlisting}
theorem twoBandParentLaurentMatrix_sq (gap hopping : ℝ) :
  twoBandParentLaurentMatrix gap hopping *
      twoBandParentLaurentMatrix gap hopping =
    twoBandLaurentRadicand gap hopping •
      (1 : Matrix (Fin 2) (Fin 2) (LaurentPolynomial ℂ))
\end{lstlisting}
The proof uses matrix extensionality, the four \texttt{Fin 2} cases,
\texttt{LaurentPolynomial.T\_add}, and normalized commutative-ring arithmetic.

\noindent
\begin{tabularx}{\textwidth}{@{}lX@{}}\toprule
Lean toolchain & Lean 4.35.0 \\
Mathlib revision & \texttt{56a005855ea17e0b0958f73ec801ae140b0a1679} \\
Additional dependencies & none \\
Principal imports & Laurent polynomials, complex basics, matrix notation \\
Validation & see \path{docs/trusted-base-audit.md} for historical exact-source and named-audit artifacts \\
Promotion & current reconciliation tree still requires promotion gates \\
\bottomrule\end{tabularx}

\section{Trusted base and axiom report}
\begin{lstlisting}
#print axioms ProjectKoios.BandGauge.Projection.twoBandParentLaurentMatrix_sq
-- [propext, Classical.choice, Quot.sound]
\end{lstlisting}
There are no repository-owned axioms or admissions in the declaration.

\section{AI, automation, and comprehension record}
The theorem selection, proof exploration, Lean tactic script, and dossier prose were
AI-assisted. Lean elaboration and kernel checking validated the exact declaration.
\texttt{ring}, \texttt{ring\_nf}, \texttt{norm\_num}, simplification, and finite-case
automation produced proof terms checked by Lean. The maintainer selected the source
claim and reviewed the displayed algebra and boundaries. No independent expert
mathematical review has occurred.

\section{Scientific and evidence boundaries}
This theorem contains no numerical input, certificate, runtime artifact, material model,
or physical validation. By itself it does not prove that the matrix is a Bloch Hamiltonian
of a material, derive its spectral projectors, or prove that a projected operator has
infinite hopping support. Those source-specific algebraic bridges are supplied separately
by \path{BG-PROJECTION-THM-003--008}.

\section{Unresolved questions}
Canonical prior-art identity remains unresolved through the References owner. Quantitative
hopping decay is outside this theorem and its companion counterexample.

\section{Review and revision history}
\begin{tabularx}{\textwidth}{@{}lllX@{}}\toprule
Date & Revision & Reviewer or process & Outcome and limitations \\
\midrule
2026-10-20 & initial & AI-assisted implementation; Lean kernel check & exact supporting
identity proved; no independent expert review \\
2026-10-20 & companion integration & Lean kernel check & full source proposition completed
separately as \path{BG-PROJECTION-THM-008} \\
\bottomrule\end{tabularx}
\end{document}
